\chapter{Deformable Bodies and Continuum Mechanics}
\label{ch:deformable}

\begin{gomdraftnotice}
This chapter is under active development and contains only a structural outline
with preliminary material. Full exposition, worked examples, proofs, and exercises
will be added in future editions.
\end{gomdraftnotice}

\epigraph{%
  When the rigid-body assumption fails, we do not abandon mechanics.
  We embrace the continuum.
}{}

\section{When Rigid-Body Assumptions Fail}

The rigid-body model fails when:
\begin{enumerate}
  \item \textbf{Internal stresses are needed}: to predict injury, fracture, or tissue damage
  \item \textbf{Deformation affects kinematics}: shaft flexibility in golf, spine compliance
  \item \textbf{Wave propagation matters}: impact biomechanics, blast injury
  \item \textbf{Fluid-structure interaction}: blood flow, synovial fluid, respiratory mechanics
\end{enumerate}

\section{Beam Models for Long Bones and Shafts}

When deformation is small and the body is slender (length $\gg$ diameter), beam
theory provides an efficient model. A schematic representation of the finite element
beam model is shown in Figure~\ref{fig:ch9:beam-fem}.

\begin{figure}[htbp]
\centering
\begin{tikzpicture}[scale=0.85, >=Stealth]
  % Continuous beam
  \draw[very thick] (0, 2) -- (6, 2.3) node[right] {Euler-Bernoulli beam};
  \filldraw (0, 2) circle (3pt);
  \node[below] at (0, 2) {Clamped};

  % Distributed load
  \draw[->, thick, red] (1.5, 2.8) -- (1.5, 2.2) node[right] {$f(x,t)$};
  \draw[->, thick, red] (3, 2.9) -- (3, 2.2);
  \draw[->, thick, red] (4.5, 2.8) -- (4.5, 2.2);

  % Finite elements
  \node at (3, 0.8) {\textbf{Discretization into FE elements}};
  \filldraw (0, 0) circle (2pt);
  \filldraw (1.5, 0.1) circle (2pt);
  \filldraw (3, 0) circle (2pt);
  \filldraw (4.5, 0.05) circle (2pt);
  \filldraw (6, 0.2) circle (2pt);

  \draw (0, 0) -- (1.5, 0.1) -- (3, 0) -- (4.5, 0.05) -- (6, 0.2);

  % Deformation
  \draw[dashed, gray] (0, 2) -- (0, 0);
  \draw[dashed, gray] (6, 2.3) -- (6, 0.2);
\end{tikzpicture}
\caption{Finite element beam model: continuous Euler-Bernoulli beam (top) and its discretization into nodal elements (bottom).}
\label{fig:ch9:beam-fem}
\end{figure}

For a Euler-Bernoulli beam:
\begin{equation}
  EI \frac{\partial^4 w}{\partial x^4} + \rho A \frac{\partial^2 w}{\partial t^2} = f(x, t),
  \label{eq:ch9:euler-bernoulli}
\end{equation}
where $E$ is the elastic modulus, $I$ is the second moment of area, $\rho$ is
density, $A$ is cross-sectional area, $w(x,t)$ is transverse displacement, and
$f$ is the distributed load.

\subsection{Modal Analysis}

The solution is expanded in modes:
\begin{equation}
  w(x, t) = \sum_{k=1}^{N} \phi_k(x) \eta_k(t),
  \label{eq:ch9:modal-expansion}
\end{equation}
where $\phi_k(x)$ are mode shapes and $\eta_k(t)$ are modal coordinates. This converts
the PDE to a system of ODEs that can be appended to the rigid-body equations of motion.

\subsection{Golf Club Shaft Flexibility}

The golf club shaft is the most accessible example of a flexible link in sport
biomechanics. A typical shaft has:
\begin{itemize}
  \item Length: $\sim$1.1~m
  \item First bending mode: 4--6~Hz (driver), 6--10~Hz (irons)
  \item Tip deflection at impact: 5--15~cm
  \item Lead/lag deflection: 2--5~cm
  \item Toe-down droop: 3--8~cm (gravity effect during downswing)
\end{itemize}

\begin{lstlisting}[language=Python, caption={Modal analysis of a golf club shaft.}]
import numpy as np
from scipy.linalg import eigh

def shaft_modal_analysis(
    n_elements: int = 20,
    length: float = 1.1,
    EI_root: float = 50.0,
    EI_tip: float = 15.0,
    rho_A: float = 0.15
) -> tuple[np.ndarray, np.ndarray]:
    """Compute natural frequencies and mode shapes of a tapered shaft.

    Uses finite element beam model with linear EI taper.

    Parameters
    ----------
    n_elements : int
        Number of finite elements.
    length : float
        Shaft length in meters.
    EI_root : float
        Flexural rigidity at root (grip end) in N*m^2.
    EI_tip : float
        Flexural rigidity at tip (head end) in N*m^2.
    rho_A : float
        Mass per unit length in kg/m.

    Returns
    -------
    tuple of np.ndarray
        (frequencies_hz, mode_shapes)
    """
    n_nodes = n_elements + 1
    n_dof = 2 * n_nodes  # displacement + rotation at each node
    h = length / n_elements

    K = np.zeros((n_dof, n_dof))
    M = np.zeros((n_dof, n_dof))

    for e in range(n_elements):
        # Local EI for this element (linear taper)
        xi = (e + 0.5) / n_elements
        EI_local = EI_root + (EI_tip - EI_root) * xi

        # Euler-Bernoulli beam element stiffness matrix
        k_e = EI_local / h**3 * np.array([
            [12,   6*h,  -12,   6*h],
            [6*h,  4*h**2, -6*h, 2*h**2],
            [-12, -6*h,   12,  -6*h],
            [6*h,  2*h**2, -6*h, 4*h**2]
        ])

        # Consistent mass matrix
        m_e = rho_A * h / 420 * np.array([
            [156,   22*h,  54,   -13*h],
            [22*h,  4*h**2, 13*h, -3*h**2],
            [54,    13*h,  156,  -22*h],
            [-13*h, -3*h**2, -22*h, 4*h**2]
        ])

        # Assembly
        dofs = [2*e, 2*e+1, 2*e+2, 2*e+3]
        for i_idx, i in enumerate(dofs):
            for j_idx, j in enumerate(dofs):
                K[i, j] += k_e[i_idx, j_idx]
                M[i, j] += m_e[i_idx, j_idx]

    # Boundary conditions: clamped at root (grip)
    free_dofs = list(range(2, n_dof))  # Remove first node
    K_free = K[np.ix_(free_dofs, free_dofs)]
    M_free = M[np.ix_(free_dofs, free_dofs)]

    # Solve eigenvalue problem
    eigenvalues, eigenvectors = eigh(K_free, M_free)

    # Convert to Hz
    frequencies = np.sqrt(np.abs(eigenvalues)) / (2 * np.pi)

    return frequencies[:6], eigenvectors[:, :6]


# Example
freqs, modes = shaft_modal_analysis()
print("Natural frequencies (Hz):")
for i, f in enumerate(freqs):
    print(f"  Mode {i+1}: {f:.1f} Hz")
\end{lstlisting}

\section{Finite Element Methods (Overview)}

For complex geometries (vertebral bodies, skull, articular surfaces), the finite
element method (FEM) provides the most general approach. While a full treatment is
beyond this text's scope, the key concept is:

The continuous domain $\Omega$ is divided into elements, and the displacement field
is approximated using shape functions:
\begin{equation}
  \bm{u}(\bm{x}, t) \approx \sum_{i} N_i(\bm{x}) \bm{u}_i(t),
  \label{eq:ch9:fem-approx}
\end{equation}
leading to the assembled system:
\begin{equation}
  M_{\text{FEM}} \ddot{\bm{u}} + C_{\text{FEM}} \dot{\bm{u}} + K_{\text{FEM}} \bm{u} = \bm{f}_{\text{ext}}.
  \label{eq:ch9:fem-system}
\end{equation}

\section*{Exercises}

\begin{enumerate}
  \item \textbf{Shaft Frequency.}
        Using the provided code, compute how the first natural frequency changes
        as the tip-to-root EI ratio varies from 0.1 to 0.9. Plot the result.

  \item \textbf{Mode Shapes.}
        Plot the first three mode shapes of the shaft. Identify nodes (zero-crossings)
        and antinodes (maximum displacement).

  \item \textbf{Timescale Separation.}
        If the golf downswing takes 0.25~s, estimate the number of oscillation cycles
        for each of the first 3 modes. Which modes are quasi-static?

  \item \textbf{(Programming.)} Add a tip mass (clubhead, 0.2~kg) to the shaft FEM
        model and recompute the natural frequencies. How does the clubhead mass affect
        the frequencies?
\end{enumerate}
